\section{Recognizing supported extensions of affine quotients}
\label{sec:period-extensions}

The following interface combines the two recognition steps. It separates a
presentation of the affine base from the presentation of a group with larger
local sections. The latter is not identified merely by displaying a finite
image of the right type.

\begin{theorem}[AA.PeriodExtension]\label{thm:period-extension}
Let $I$ be finite and let $U=\langle r_i:i\in I\rangle$ be a finitely
presented marked group with an affine marked image
$E\leq\mathbb Z^d\rtimes Q$, where the finite linear images agree. Suppose
the actual kernel of $U\to Q$ has a presentation identifying it with
$\mathbb Z^d$, and that the translation kernel $L$ of $E$ has rank $d$.
Let $M\leq L$ be $Q$-invariant, with an explicitly given finite lattice
basis represented by words in $U$.

Suppose the equations of AA.ParabolicPeriods hold for $M$. Put $B=E/M$.
Choose a finite antichain $\mathcal F$ of label supports. For each
$S\in\mathcal F$, suppose there is a deletion $p_S:B\to B_S$ and a
marked epimorphism $q_S:V_S\to B_S$ from a group generated by the
markers on $S$, injective on every proper marked
subgroup of $V_S$. Write $N_S=\ker q_S$.

Let $P$ be a marked presentation with the following properties:
\begin{enumerate}
\item Its relators hold in the compatible product of $B$ and the $V_S$
with respect to $p_S,q_S$.
\item Every local presentation for $V_S$ holds in $P$.
\item The image in $P$ of each $N_S$ is normal in $P$.
\item Killing these local kernel images leaves exactly the presentation
of $U$ with the basis words of $M$ set equal to one.
\end{enumerate}
Then $P$ is marked-isomorphic to that full compatible product. It has IP,
and its parabolic projection kernels are the coordinate products
\[
 \ker(P_J\to B_J)=\prod_{S\subseteq J}N_S.
\]
Each local subgroup $P_S$ is exactly $V_S$. If $[L:M]$ and all $N_S$ are
finite, then
\[
 |P|=|Q|[L:M]\prod_{S\in\mathcal F}|N_S|,
 \qquad |P_J|=|B_J|\prod_{S\subseteq J}|N_S|.
\]
The geometric and deletion conclusions hold precisely under the additional
hypotheses stated in UA.SK.Geometry and UA.SK.Delete, respectively.
\end{theorem}
\begin{proof}
AL.2 identifies $U$ with $E$ using the actual free abelian kernel
presentation and its surjection onto the rank-$d$ translation lattice.
The normal closure of the chosen basis translations of $M$ in $E$ is
$M$: it contains their additive span, and invariance under the linear
image shows that this span is already normal. Thus adjoining these words
identifies the whole base presentation with $B$, not only with a quotient
of $B$. AA.ParabolicPeriods gives IP in $B$.

The stipulated deletions and proper-parabolic injectivity are exactly the
hypotheses of UA.SK. That theorem constructs and generates the full
compatible product, identifies every projection kernel, and proves IP.
The four presentation hypotheses are exactly UA.Recognition. Its proof uses
normality and the faithful separate local coordinates to exclude every
remaining kernel; it proves $P$ is that compatible product. The antichain
clause UA.SK.Local identifies its local subgroups. When the displayed
indices are finite, $|B|=|Q|[L:M]$, and UA.SK gives the stated orders.
The last sentence invokes the additional proved hypotheses of the two
named corollaries; IP alone does not establish a hypertope's flag property.
\end{proof}

The theorem also applies with no local supports, in which case $P=B$.
For an affine base with an elementary normal-form proof, that proof can
replace the Schreier computation in the first paragraph: the required input
is the exact marked isomorphism $U=E$, not a particular algorithm for proving
it. This is how the square groups enter the two-power applications. The
cubic $\{6,4\}$ and $\{6,6\}$ bases instead retain their two full exact
kernel-presentation programs. Their lattice and overlap programs verify the
separate arithmetic hypotheses. No one of these certificates replaces the
others.

This theorem assembles already proved normal-kernel and supported-extension
arguments with the period criterion. The eight applications below retain
their source authorship. The high-type and order-1024 classifications also
need the forest, cyclic-core and exhaustive-kernel proofs; they are not
consequences of this construction theorem alone.
